\documentclass[oxford]{ahsan}
\title{Recovery From Transient Failures Of The Apollo Guidance Computer}

\begin{document}

\maketitle

\section{ABSTRACT}

In the Apollo Guidance Computer, nearly one hundred thousand word transfers occur each
second. A random error rate of one in \(10^{12}\) actions would be considered good in today’s
technology, but at that rate an error might occur within several hundred hours. Added
to the random error rate are externally induced errors, including power and signal
interface transients, program overloads, and operator errors. The incidence of induced
errors has so far been very much greater than random errors (if indeed there have been
any at all) in AGC experience.

The AGC is a control computer, and has been designed to detect and to recover from
random or induced transient failures. The techniques used are the subject of this
paper.

The paper outlines the character of the computer and its system software to a depth
sufficient for the discussion. The built-in malfunction alarm logic is discussed,
along with the software- based computer Self Check. These systems, upon detection of a
failure, force an involuntary re- ordering of the signal interface, and an involuntary
transfer to the restart program logic. From this point, the software, without recourse
to suspect information in the central processor, reconstructs the output interface
conditions, and reconstitutes the control processes in progress at the time of failure.
There are a number of ground rules which limit the recovery capability, based on
presumptions about the nature of the failure, These are presented.

The software associated with a restart is described, with a typical program flow
derived from an Apollo Mission Program. The amount of memory assigned to restart
protection is stated. Several interesting sidelights are briefly discussed such as
manual break-in to prevent restart looping, the adoption of the restart technique for
scheduling the termination of active programs, and the use of the failure recovery
technique to remove a temporary computer overload.

\section{Introduction}

The Apollo Guidance Computer (AGC) is part of the Guidance, Navigation and Control
System of the Apollo Spacecraft. There are two systems, one in the Command Module and
one in the Lunar Module. This paper is about the restart feature of the hardware and
software. The intent of the paper is to explain what restart protection is, why it is
done, what things might have been done instead and implications, large and small, of
various design decisions.

The name restart refers to a built-in recovery procedure which responds to hardware
detected computer malfunctions. Basically, the hardware transfers control to a specific
location in the AGC memory. The software then restarts the programs in progress by
reference to stored data which points to the last restart point passed by each program
before the malfunction was detected.


\section{The Apollo Guidance Computer}

The AGC is a control computer. It receives incremental angular and velocity information
from an Inertial Measurement Unit, from a telescope, a sextant, and in the case of the
Lunar Module it receives range, range rate and angular data from a Rendezvous Radar and
a Landing Radar. There is a Display and Keyboard (DSKY) which is illustrated in Figure
1. There is an uptelemetry receiving system, and a downtelemetry system.

There are a number of discrete signals sent and received, indicating, for example, such
events as liftoff, accelerometer failure, a-rd  vehicle stage separation.

The computer directly controls the numerical displays, the rocket engine gimbals, the
attitude jet commands, angular commands to the radar, optics, and inertial measurement
unit, attitude error displays, engine throttle and engine on/off commands, among other
things.

The computer performs many functions, almost all of which are variations on four
purposes:
\begin{enumerate}
    \item to maintain and improve knowledge of position and velocity,
    \item to compute changes in positions an velocity by thrust vector control to
        carry out various mission phases,
    \item to control vehicle attitude, and,
    \item to compute and display a wide variety of information necessary for astronaut
        control.
\end{enumerate}

All of these functions are under the control of the astronaut, through the DSKY.
There are about 40 separate programs available to him. These programs are
function oriented, and are strung together by the astronaut according to his
needs by a succession of keyboard actions. In addition to the programs, there
are approximately 30 routines which can be run simultaneously with the programs
to yield information as desired. In general, only one program can occupy the
machine at one time. There is an exception to this which will be discussed in
detail later because it is relevant to the restart program.

The progress through a program or a routine is characterized in most cases by a
sequence of displays. There are many decision points, at which the relevant data is
displayed and a flashing announcement indicates that an astronaut response is required.
The usual astronaut response is to Proceed with the computational path and values
displayed.

Provision is made to put programs “to sleep” until a desired response or other
necessary external event occurs. The display system is a special and intricate
part of the AGC software, since display information, possibly requiring an
astronaut res- ponse, may be “buried” under display information generated by an
auxiliary routine, which in turn may have been interrupted by a priority
display, or by astronaut keyboard activity. There is aKey  Release light which
informs the astronaut about buried information.


\section{AGC Memory}

The AGC memory is composed of three significantly different components. The first
consists of 36,864 words of non-alterable program storage, called fixed memory. The
contents of the memory locations are fixed at manufacture time. The question of content
is one of physical threading or bypassing of wires through a sequence of tape wound
magnetic cores. Restart protection presumes that fixed memory is intact.

There are 2040 words of program-alterable memory, called erasable memory. These consist
of a conventional destructive read out ferrite core matrix. The question of content is
one of the polarity of magnetic fluxset  in theindividual core. Restart protection
presumes that erasable memory is intact, although a small portion is redundantly stored
and compared during a restart.

There are seven Central registers, designed for high speed operation. The names of
these registers are in Table 1. These registers are composed of flip flops, formed by
interconnecting logic NOR gates, with content dependent on continuity of electrical
signal from the time thememoryis set, to the time that it is read. Restart protection
presumes that the content of these registers may have been destroyed

There are sixteen input/output channels, similar to the central registers, employing
logic flip flops.

\section{The AGC Software Operating System}

There is one central processor. There is an Executive program which schedules Jobs,
according to program assigned priority. Each Job is required to periodically
interrogate the NEW JOB register which indicates the presence of Jobs of higher
priority. If there is a Job of higher priority awaiting execution, the current Job is
replaced. At the completion of a Job, Jobs of successively lower priority are executed
until finally the Job of lowest priority called DUMMY JOB is continuously executed. A
computer Self Test program, designed to test the reading of fixed memory and the
reading and writing of erasable memory can be executed during this time. When a Job
request is made, a fixed length set of erasable core is assigned to that Job and is
kept intact until the Job is terminated. Once a Job is active, it is not possible to
recognize it among the other Jobs in various stages of completion in the Executive.
This is a difficulty when, as is frequently necessary, a subset of Jobs active in the
Executive must be terminated.


\section{Interrupt Mode}

There are several events, such as a keyboard depression, the overflow of various timing
counters, the availability of data at the radar interface, etc., which cause an
involuntary transfer of control to specific locations in the AGC. The computer is then
operating in interrupt mode. After storing away the information needed to continue the
interrupted program, the interrupt program either can complete a short computational
task before exiting from the interrupt mode, or it can set up a Job to be done by the
Executive after terminating the interrupt. There are several timing registers which can
be set under program control, and which cause interrupts when they come  due. These are
used to schedule time critical events. An example is TIME 3: when the TIME 3 register
overflows, it means that a so-called Task must be executed. This Task is performed
under interrupt. The Waitlist program isused  toset  the TIME 3 counter and to maintain
a list of starting addresses for these Tasks. A typical example is the READACCS Task
during powered flight. This Task reads the accelerometer velocity increment registers,
makes a Job request to the Executive which will complete a rather lengthy series of
computations, makes a request to the Waitlist to reschedule itself for 2 seconds from
now and then terminates the interrupt mode. When this happens, the Executive resumes
where it left off, unless another interrupt has occurred in the meantime. There is an
instruction, INHINT, that delays interrupts until execution of another instruction,
RELINT.

A restart may be considered to be a special case of interrupt. However this and only
this event will interrupt a program already operating in interrupt mode, or a program
operating with interrupt inhibited. Also, restart does not return control to the
interrupted program.

\section{Need For Restart}

In the AGC, nearly one hundred thousand word transfers are performed each second. A
random error rate of one in 10 12 actions would be considered good in today’s
technology, but at that rate, an error might occur within several hundred hours. Added
to the random error rate are externally induced er- rors, particularly power
transients. The possibility of error presents a requirement for a means of recovery,
for the AGC itself would still be perfectly useful given that it could be re-initiated
into the program. To meet the need for re-initiation, the AGC is equipped with a
restart feature comprising alarms to detect malfunction and a standard initiation
sequence which leads back into the programs in progress.

\section{Logic Associated With Restart}

A restart is caused by a detected computer malfunction. When a restart occurs there is
an immediate involuntary transfer of control to location 4000, in fixed memory. Also
several signals internal to the AGC are set to insure a known logical state. All the
output channels, which include the following functions, are reset (cleared):

\begin{itemize}
    \item individual jet commands 
    \item display relay commands 
    \item engine off and engine on commmds 
    \item various warning and status lights 
    \item various moding commands to other spacecraft systems 
    \item engine gimbal motion commands 
    \item down telemetry word in process of transmission
\end{itemize}

Clearing these output channels at least places the interface in a known condition.
Otherwise, because of the restart, it is not certain what their status would be. It is
left to the software to reconstruct the correct output interface as rapidly as
possible.

The coding following the transfer to location 4000 is designed to reinitialize the
system programs. Incremental clocks are set to get rapid interrupts to start the
various time dependent functions. Jobs in the Executive and Tasks in the Waitlist  are
cleared out. The display panel is blanked except for the restart light. The various
discrete monitors are started. The digital autopilot is initialized. The desired state
of the engine is checked and the output bit is properly set.

Following this sequence, the program tests a phase pointer for each of five restart
groups. Since there may be several Jobs in progress, running without synchronization,
it is necessary to be able to advance restart points in different programs
independently. To handle this, there are five independent restart groups, each with its
own phase pointer. Jobs and Tasks that may run simultaneously must be in different
groups. The pointers are stored redundantly. If each phase pointer passes a redundancy
check, the pointers are used to restart the groups.


MAJOR TASKS FOLLOWING A RESTART: 
\begin{itemize}
    \item Save diagnostic informat ion, 
    \item Make executive and waitlist dormant, 
    \item Initialize the system program, 
    \item Check for lock out, 
    \item Test erasable, 
    \item Determine if programs were in process and where to re-enter them.
\end{itemize}

\section{Malfunction Detection Devices}

There are six malfunction detection devices that cause a restart. They are discussed
below.

Parity failure detection is the most powerful. A single odd-parity bit is added to each
fixedmemory word at manufacture time, and to each erasable word at write time. Parity
is tested at read time. If the test is failed, the restart sequence begins.

The Night Watchman failure detection works as follows: the NEW JOB register, mentioned
earlier, must be periodically tested by any program which is following the rules which
make the Executive work. This register is “wired” and if it is not tested often enough,
the Night Watchman circuit causes a restart.

The TC Trap detects the endless one- instruction loop caused by transferring control to
a location L which contains the instruction “transfer control to location L”.

The Oscillator Fail is caused by a stopping of the timing oscillator.

The Voltage Fail circuits monitor the 28-, 14-, and 4-volt power levels which drive the
computer.

Rapt Lock detects excessive time spent in interrupt mode, or too much time spent
between interrupts.


\section{Assignment of Restart Points}

At the time of the detected malfunction, the computer may or may not be in interrupt
mode. In any case, there was a list of Jobs and a list of Tasks, with associated
incremental times, awaiting execution. An uncompleted interrupt mode program is
usually not restarted, with the exception of Tasks and certain autopilot functions.
This means that a downtelemetry word may be garbled, or a key depression may have to be
repeated. The restart pointers allow the restart program to pick up proper Job and Task
information for a proper restart of the active programs.

Lists of such restart information are stored in fixed memory.’ The lists are called
Restart Tables. During the execution of the program, a pointer is advanced to
successively point at the proper restart information for that phase of the solution.
This is part of the design of the program.

As an alternative to setting a pointer to pick up the information from the proper point
in the phase table, other options are available. Often, the instruction that began the
last display sequence is an acceptable restart point. This is also a convenient place
since the same information relevant to restarting the program must be saved by the
Display Program for its own purposes. This type of restart is frequently used and has a
special defining format. Another frequently used technique is to define the restart
point to be the instruction to be executed next. This does not require use of the
Restart Table.

It is occasionally necessary to simultaneously restart combinations of Jobs and Tasks.
In other circumstances it is necessary to specify the restart information indirectly,
that is, the restart pointer picks up the address where restart information can be
found instead of picking up the information directly.

The information for defining a restart point is packed in binary statements which are
decoded by a subroutine. The subroutine advances the pointer or performs whatever other
operations are needed to make provision for a possible restart. Thus, between
invocation of this subroutine, the in- formation about how to restart the software
waits, stored in erasable locations, until the mainline programs again use the
subroutine to advance the restart information.

\section{Testing for Restart Protection}

Restart logic is tested during the development of a flight program by causing restarts
to happen at specified times during simulations of mission sequences. These Situations
are at the instruction level of the AGC and provide detailed checks on program design.
Restarts can also be caused during simulation using a real AGC connected to test
equipment. These restarts are benign, in the sense that there is no malfunction, and
therefore no possibility of destruction of memory between malfunction and detection.
There is no ability at present to insert the restart between specified instructions,
although this would be a useful test option.

There has been some discussion and preliminary planning of a “Restart Analyzer” which
examines the code at simulation time to detect and complain when the execution of the
instruction violates restartability rules. There are several problems to implementing
such an analyzer; first, the statement of a set of sufficient conditions has not yet
been achieved; secondly, the analyzer would greatly increase the run time of a
simulation, although it would probably decrease total test time.


\section{Application  of Restart Techniques for the Control of Program Sequencing}

Rendezvous Navigation uses measured data to update knowledge of relative position and
velocity. The information is a combination of angle, range and range rate information
from various devices on the spacecraft. This program in the AGC is called P20.

Rendezvous Targeting uses various algorithms to compute velocity changes required to
complete the rendezvous. There are several targeting programs, which correspond to
different techniques and phases of rendezvous. In all there are six astronaut-callable
programs that do Rendezvous Targeting.

During the development of the specifications for these programs, it became apparent
that the time spent in the targeting programs should also bused for navigation
measurements, Thus a requirement existed for simultaneous running of P20  and  one of
the various targeting programs, say P34 (Transfer Phase Initiation). It soon became
apparent that it was important to be able to start and stop navigation and targeting
functions independent of each other.

The mechanization of this simultaneous pro- gram capability, including the ability to
selectively terminate part of the functions, is similar to the procedures used in
recovering from a detected malfunction. When any program is to be terminated, all
future Tasks yet to go active are cleared out of the Waitlist. All Jobs whether they
are currently active, sleeping, or waiting for their priority to come to the top, are
terminated. Then, the restart group or groups associated with the functions that are to
be continued are restarted. That is, they are activated according to the Restart
pointer associated with that group.

Another example of the application of this technique occurs during powered flight
where, in at least one non-predictable circumstance, all functions, no matter what
their stage of completion, have to be terminated, except for the integration of
position and velocity. This is easily implemented by terminating all computations and
restarting the restart group associated with integration.

Thus the logic associated with restart protection serves a significant function
independent of the question of attempting to recover from detected hardware failures.
There are, of course, other design paths that could have been followed to obtain the
ability to selectively terminate active programs. It is interesting to note that the
Executive system, which does not maintain a record of starting ad- dresses (or other
clues as to the origin of the Job) beyond the time when a Job is first started,
probably strongly influenced the decision to selectively terminate Jobs by the restart
technique.


\section{Software Abort Conditions}

Programming dead ends, which are not expected to be encountered in a checked-out
program, under proper use, are broken into two categories; the first and largest
category involves conditions which probably will recur again and again if the program
is restarted. These conditions are handled by abandoning the program in progress and
flashing a request to select another program. Also an alarm light is lit and, by
keyboard action, the astronaut can find ,the code number of the alarm that caused the
program abort.

There is a small class of problems centering around overloading the Executive or
Waitlist, where it is likely that trying the program again via the restart logic will
be successful. This is done. A large class of auxiliary programs called Extended Verb
Routines can, through astronaut keyboard command, be run simultaneously with a program.
If an Executive overload occurs, and if the condition is partially the result of a
so-called Extended Verb Routine, restart will solve the problem because, by design,
extended verbs are not restartable. However, to differentiate the procedure from
hardware malfunctions, the restart warning light is not lit.

\section{Diagnostic Information }

The contents of the program counter and the bank registers, which together point to the
instruction that was being executed when a restart occurred, are saved by software
means. They may, however, have been destroyed by the malfunction. 

It is not possible to differentiate between the various malfunction detectors while the
computer is in a flight configuration. There is information brought out through a test
connector which permits isolation to a particular failure detection mechanism during
ground test.

There's,  of course, the possibility that restarts are self-generating. If the same
problem still exists on re-executing the code, there is a potential endless loop.
In view  of this, it would probably be better to store away the diagnostic information
only once, on the first pass through the restart logic. There is a register which is
incremented on each pass through the restart logic.


\section{ Protection Against Endless Loops}

Since the restart logic tries to do a piece of program over again, the probability of
endless looping must be dealt with. When this occurs, the loop can be so tight that
ordinary means of control- ling the computer are not available. For example, the
Display Keyboard, or the telemetry input programs, are not serviced. Even turning off
the computer does not necessarily break the loop. Therefore, the first thing that is
done in the restart logic is to check for a simultaneous depression of two buttons
(input keys). If the user depresses these keys, the logic diverts the program to an
inert condition, called Fresh Start.

\section{Self Test versus the Restart Philosophy}

Self Test tests the ability to read and write into erasable locations. The process
consists of storing the contents of an erasable location in another location, writing
and reading all zeros, then, writing and reading all ones, then retrieving the original
contents and writing back into the register.

Self Test presents an interesting dilemma when juxtaposed against the philosophy of
restart protection. There are two reads and two writes (the ones and zeros) that if
unsuccessfully completed will be detected. On the other hand, there are two reads and
two writes that must be successfully completed at each test in order not to destroy the
contents of the register. Self Test cannot itself detect an error in this part of its
operation. If a transient failure occurs while Self Check is testing erasable, there is
as much chance that the contents of a register will be destroyed as there is that a
transient failure will be detected and erasable memory preserved. If there has been a
hard failure, self test will detect it, and the contents of the register are not as
important as finding the failure. Parity detection enters into this discussion, since
at least half, and probably much more than half, of all read and/or write errors will
cause parity failure, which will itself cause a restart. This does not alter the
fundamental point that if transient failures are significant enough to necessitate
designing restart logic, then running Self Test has a significant chance of destroying
erasable information. Thus searching for hard failures by testing read and write logic
may cause undetected destruction of memory due to transient failures. Probably the best
thing to do is to do Self Test before beginning significant mission phases, but not
otherwise.

\section{Will Restart Work?}

If a computer failure is such that erasable memory information is still intact, and if
the failure is indeed transient, then there is great confidence that the restart logic
will bring the program through the transient.

Will the failure mechanism be so generous? We do not know. No restart has occurred in
flight. Those that have occurred in test have occurred with non-rigorously checked-out
test programs and are usually traced to software-invoked TC TRAP failures. This was the
method of dealing with program “dead ends” and is no longer used. Little definite
information is available.

We can gain some feeling for the problem by considering the mechanism used to detect
the malfunction.

If there is a parity failure on reading fixed memory, the chances of recovery are very
good. This is because the failure would not destroy the memory, and a subsequent
attempt will be successful if the failure mechanism has disappeared. More important,
the failure is detected immediately, before other damage is done.

A parity failure on reading an erasable location has far less chance of successful
restart because the contents of the erasable is suspect, since the memory is a
destructive read system, and the register is refilled with the word that just failed
parity with the parity bit corrected.

If the restart is caused by detecting a TC TRAP or a Night Watchman alarm, it is likely
that the actual failure has occurred many, many, many instructions ago, and that the
program has been running wild for a considerable time. Under these circumstances it is
not likely that erasable memory is intact.

If the restart is caused by a detected voltage dropout, there is enough capacitance in
the power circuits to keep the central registers intact until the completion of the
current memory cycle. Therefore, transfers of information to erasable memory will be
completed with good information. Intentional failures of power levels have been tried
with successful results, and unplanned voltage dropout is the failure most often
experienced in ground checkout. Of course, the oscillator also stops, so that there is
an unknown time mismatch between the AGC clocks and real time. There is no way to
correct this within the ground rules of restart, but ground tracking can detect and
correct this loss of time synchronization. Thus, the chances of a successful restart
vary greatly with the source of the detected failure. There are two failure mechanisms
that seem to offer the best chance of successful restart, It is important to note that
these two failures could be handled by hardware circuits as an alternative to the
approach described in this paper. First, parity failure on reading fixed memory. This
could be dealt with by reading the location again, which would be much simpler. Second,
voltage transients could be solved by saving flip flop information, then letting the
AGC start up where it left off.

\section{What  Should the Astronaut Do if the Restart Light Comes On?}

If a restart occurs in flight, the astronaut and the ground will be warned by a restart
light on the DSKY. Furthermore, the G\&N-Caution Light will come on, and if a sequence
of restarts occur in a short interval, the AGC Warning Light will come on, along with a
light called the Master Warning Light. The program may or may not give evidence of
memory loss. If it does, then erasable memory should be reinitialized with help from the
ground. If the program continues to work well, erasable memory should still be
initialized at the earliest convenient time. Self Test should be run, and all other
means of diagnosis should be called upon. The computer should not be presumed defective
although, depending on the flight circumstances, the best strategy might be to go to
backup equipment until confidence in the computer can be reestablished.

\section{The Cost of Restart Protection}

The principle involved in the development of this program is simple. If a detected
hardware malfunction occurs, then, since it is highly likely that the failure is
transient, it is worthwhile to try to continue the computation. The alternative is to
abandon the computer until a new erasable load is transmitted to the ground. The
specific techniques developed are complicated and increase the difficulty in
delivering checked-out flight software. Some of these effects can be described
quantitatively First, the programs associated with restart occupy 1048 locations of
the 36,864 words of fixed memory.

Programs that are restart protected must contain phase *update  instructions. These
occupy 334 locations in fixed memory. Thus about 4\% of the memory budget is assigned to
restart protection. There is an indirect memory cost associated with designing a
program to be restartable.

In terms of testing, some 40\% of program testing and 1\% of mission phase testing
includes restart testing.

There is another cost. Restart protection requires a knowledge of the overall working
of the machine, and an intimate knowledge of areas of program that would otherwise not
be of significance in the development of a particular software program region.


\section{Conclusion}

The establishment of restart protection in a computer such as  the AGC presents somewhat
of an enigma. In a total of over 25 hours of space flight, the computer has yet to have
a transient failure from which the restart feature could be called on to demonstrate
its worth. This could well be the experience for the whole  Apollo program. We have seen
that the provision for restart in the computer program complicates the generation and
test of program. We have seen that there is a significant class of transient failure
events which restart will probably fail to cure. And yet only one successful recovery
by restart might save a mission.


\end{document}
